
\documentclass{article}
%%%%%%%%%%%%%%%%%%%%%%%%%%%%%%%%%%%%%%%%%%%%%%%%%%%%%%%%%%%%%%%%%%%%%%%%%%%%%%%%%%%%%%%%%%%%%%%%%%%%%%%%%%%%%%%%%%%%%%%%%%%%
\usepackage{amsmath}
\usepackage{harvard}

%TCIDATA{OutputFilter=LATEX.DLL}
%TCIDATA{Created=Sunday, March 24, 2002 09:08:29}
%TCIDATA{LastRevised=Sunday, March 24, 2002 09:09:03}
%TCIDATA{<META NAME="GraphicsSave" CONTENT="32">}
%TCIDATA{<META NAME="DocumentShell" CONTENT="Articles\SW\Standard LaTeX Article (Harvard)">}
%TCIDATA{CSTFile=LaTeX article (bright).cst}

\newtheorem{theorem}{Theorem}
\newtheorem{acknowledgement}[theorem]{Acknowledgement}
\newtheorem{algorithm}[theorem]{Algorithm}
\newtheorem{axiom}[theorem]{Axiom}
\newtheorem{case}[theorem]{Case}
\newtheorem{claim}[theorem]{Claim}
\newtheorem{conclusion}[theorem]{Conclusion}
\newtheorem{condition}[theorem]{Condition}
\newtheorem{conjecture}[theorem]{Conjecture}
\newtheorem{corollary}[theorem]{Corollary}
\newtheorem{criterion}[theorem]{Criterion}
\newtheorem{definition}[theorem]{Definition}
\newtheorem{example}[theorem]{Example}
\newtheorem{exercise}[theorem]{Exercise}
\newtheorem{lemma}[theorem]{Lemma}
\newtheorem{notation}[theorem]{Notation}
\newtheorem{problem}[theorem]{Problem}
\newtheorem{proposition}[theorem]{Proposition}
\newtheorem{remark}[theorem]{Remark}
\newtheorem{solution}[theorem]{Solution}
\newtheorem{summary}[theorem]{Summary}
\newenvironment{proof}[1][Proof]{\textbf{#1.} }{\ \rule{0.5em}{0.5em}}
\input{tcilatex}

\begin{document}

\title{The Title of a Standard LaTeX Article}
\author{A. U. Thor \\
%EndAName
The University of Stewart Island}
\maketitle

\begin{abstract}
We study the effects of warm water on the local penguin population. The
major finding is that it is extremely difficult to induce penguins to drink
warm water. The success factor is approximately $-e^{-i\pi }-1$.
\end{abstract}

\section{The Harvard Bibliography System}

\subsection{About This Shell}

This shell document provides a sample layout of a Standard LaTeX Article.
The Typeset Bibliography Choice has been set to BibTeX and the Harvard
package has been added. The documentation for this package is \qExtProgCall{%
DVI Previewer}{harvard.dvi}{TeX Documentation}{%
latex/contrib/supported/harvard/harvard.dvi}{}{} (Ctrl-Click the preceding
link to view the documentation).

Changes to the typeset format of this shell and its associated \LaTeX{}
formatting files (\texttt{article.cls }and \texttt{harvard.sty}) are not
supported by MacKichan Software, Inc. If you wish to make such changes,
please consult the \LaTeX{} manuals or a local \LaTeX{} expert.

If you modify this document and save it as ``Standard LaTeX Article
(Harvard).shl'' in the \texttt{Shells\TEXTsymbol{\backslash}Articles%
\TEXTsymbol{\backslash}SW} directory, it will become your new Standard LaTeX
Article (Harvard) shell.

\subsection{Compiling a Bibliography}

\begin{enumerate}
\item Save this document using a name of your choosing.

\item Choose \textsf{Typeset, Compile.}

\item Check \textsf{Generate a Bibliography.}

\item Choose \textsf{OK.}

\item Choose \textsf{Typeset, Preview.}
\end{enumerate}

\subsection{Citations and the Bibliography}

The standard citation mechanism that is supported by the SW \textsf{Insert,
Field, Citation} mechanism is the parenthetical citation form, using the
LaTeX macro \verb|\cite| and is as follows:

As \cite{agsm} and \cite[Annex~B]{latex:guide} describe . . .

\subsubsection{\texttt{\TEXTsymbol{\backslash}citeasnout}}

The Harvard package also supports using a citation as a noun using the LaTeX
macro \verb|\citeasnoun|. These mechanisms are not supported directly
through the SW interface. They are entered as TeX fields. Here are some
examples that you can copy and modify.

As \citeasnoun{agsm} and \citeasnoun[Annex~B]{latex} describe . . .

\subsubsection{\texttt{\TEXTsymbol{\backslash}possessivecite}}

You can use the citation as a possessive noun phrase with the LaTeX macro 
\verb|\possessivecite|. Multiple citations are not permitted.

\possessivecite{latex} description of this feature is . . .

\subsubsection{\texttt{\TEXTsymbol{\backslash}citeaffixed}}

You can use the LaTeX macro \verb|\citeaffixed| command to allow text to be
affixed inside the beginning of the parenthesis of a parenthetical citation.

BibTeX manuals \citeaffixed{latex,agsm}{e.g.} describe . . .

\subsubsection{Changing the default citation command}

As previously mentioned, the standard citation mechanism supported by SW
used the LaTeX macro \verb|\cite|. If you prefer to use the citation as a
noun or the citation as a possessive noun phrase throughout your \emph{entire%
} document, you can change the default SW behavior. \ This is done by adding
an appropriate macro to the preamble of your document. \ For example, if you
wish to use citation as a noun throughout your document, select \textsf{%
Typeset, Preamble }and add a line containing \verb|\let\cite=\citeasnoun|.
Be careful when editing the preamble (see the warning in the LaTeX Preamble
dialog box). Use \verb|\let\cite=\possessivecite| in the document preamble
if you want all citations to appear as a possessive noun phrase. With one of
the macros in the document preamble, simply use the normal SW citation
mechanism and the citations will appear as nouns or possessive noun phrases
when you compile your document and generate a bibliography.

\bibliographystyle{kluwer}
\bibliography{harvard}

\end{document}
